\section{Discussion}
\label{sec:discussion}

This tutorial argued that the SE community should adopt the modern causal inference toolkit as the standard for intervention-outcome questions involving observational data.
We now reflect on what the worked example reveals about current practice, discuss implications for different audiences, and acknowledge limitations and open frontiers.

\subsection{Lessons from the Worked Example}

The progression of estimates in Table~\ref{tab:all-estimates} is the paper's most concrete argument for why identification strategy matters.
The na\"ive cross-sectional comparison yields $+911\%$; adding snapshot controls reduces this to $+158\%$; conditioning only on DAG-justified pre-treatment covariates further reduces it to $+101\%$; and the staggered DiD estimates land at $+32\%$ to $+52\%$.
Each step removes a specific, named source of bias---temporal ambiguity, collider conditioning, time-invariant confounding, secular trends, common shocks---and each reduction is substantial.
The design, not the estimator, drives the change: Within the cross-sectional family, switching from OLS to PSM to IPW may improve the estimate but they are all plagued by the same conditional ignorability assumption violated by unmeasured confounders.
The magnitude of this shrinkage matters beyond methodological tidiness.
For many intervention-appraisal questions---how to allocate resources, whether to invest in a technology, whether to adopt a practice---the decision-relevant quantity is \emph{how much}, not merely the direction of the effect~\cite{lundberg2021estimand}, and an order-of-magnitude-inflated estimate can lead to qualitatively different conclusions even when the sign is preserved.

This pattern is unlikely to be unique to our setting.
Many MSR studies query the GitHub API once, obtain a cross-sectional snapshot, regress an outcome on a treatment indicator with available controls, and report the surviving coefficient as suggestive of an effect.
The temporal DAG analysis (Section~\ref{sec:temporal-collapse}) revealed that this workflow destroys the temporal ordering required for principled covariate selection: Variables that are confounders in their pre-treatment incarnation become colliders in their post-treatment incarnation, and a single snapshot conflates the two.
A practical recommendation follows: MSR researchers should invest in constructing longitudinal datasets---mining git histories, leveraging GHArchive time series, recording adoption timing---rather than relying on atemporal snapshots.
The marginal cost of building a panel is modest; the marginal gain in identification is large.

The worked example also produces a substantive finding.
Under the staggered DiD design, systematic AI coding tool adoption---as signaled by AI configuration files committed to the repository---is associated with a 32--52\% increase in monthly commits (ATT).
This estimates the effect of the \emph{entire bundle}: The AI tools themselves, the workflow reorganization that accompanies structured adoption, and the team characteristics that lead to early integration (Section~\ref{sec:study-setting}).
It does not isolate the effect of AI code generation alone.
Pre-treatment trend tests support the parallel trends assumption, but the estimate is likely conservative due to control-group contamination: Individual developers using AI tools without project-level configuration files are unobservable, blurring the treated--control contrast and attenuating the estimate toward zero.
We offer this as a credible first estimate that future work can refine with richer treatment definitions, complementary identification strategies, or alternative outcome measures.
The $+32\%$ to $+52\%$ range is also consistent with the field-experimental evidence base on AI coding tools (Section~\ref{sec:triangulation}), most directly with the $+26\%$ pooled estimate from the three firm-level randomized trials in \citet{cui2025effects}, and sits above the null-to-slowdown findings of \citet{becker2025measuring} on experienced maintainers of mature open-source repositories.
The na\"ive cross-sectional $+911\%$ is inconsistent with every published RCT by roughly an order of magnitude, illustrating that the design-based estimates are not only methodologically defensible but also the only estimates in this paper that triangulate with independent experimental evidence.

\subsection{Implications}

\paragraph{For Reviewers and Editors.}
When a paper poses an intervention-outcome research question, reviewers should require an explicit discussion of the extent to which the evidence supports a causal claim and the assumptions under which it does so.
A regression coefficient presented as ``the effect of $X$ on $Y$'' without an identification argument is not a causal finding---it is a descriptive association whose causal content is unknown.
The empirical standard in Appendix~\ref{sec:appendix-standard} provides a concrete checklist: Does the study state its causal question? Articulate a target trial? Define the estimand? Name the identification strategy and its assumptions? Conduct sensitivity analysis? Engage with alternative explanations?
We encourage editorial boards of SE venues to formally adopt this standard---or one like it---into the ACM SIGSOFT Empirical Standards collection, following the precedent of editor-endorsed guidance in respiratory and critical care medicine~\cite{lederer2019control} and emerging norms in psychology~\cite{grosz2020taboo} and epidemiology~\cite{hernan2018cword}.
This does not mean that every empirical SE paper must become a causal study; descriptive, predictive, and qualitative studies remain indispensable.
The standard applies specifically when a study asks a causal question, and a lot more empirical SE papers can possibly reframe their research question to be causal, as many of them already have a causal motivation in mind (see Section~\ref{sec:bg-adoption}).

\paragraph{For Researchers.}
The causal inference toolkit can seem daunting, but the barrier to entry is lower than it appears.
The pragmatic stance in Section~\ref{sec:guide} provides a workflow; the FAQ in Appendix~\ref{sec:appendix-faq} addresses common objections and recommends software tools; and a new generation of accessible textbooks~\cite{cunningham2021causal, huntingtonklein2021effect, hernan2020causal, angrist2009mostly} makes self-study feasible.
The most important first step is not mastering a specific method but adopting a way of thinking: Before fitting a model, articulate the target trial, sketch a DAG, ask what variation in the data identifies the causal effect, and reason whether the identification assumption would be plausible in the specific study setting.
This exercise alone---even without changing the statistical method---often reveals problems invisible to the standard regression workflow: a key covariate is a mediator, the treatment is poorly defined, or the research question conflates two distinct estimands.
Causal inquiry is inherently iterative: The first identification strategy will have gaps, and revision is the norm, not the exception.

\subsection{SE's Advantages for Causal Inference}

Our message is not that SE is behind---it is that SE has untapped potential.
Several structural features of SE data make it unusually well suited to the design-based methods this tutorial advocates.
First, SE data is inherently \emph{panel data}: Repositories have commit histories spanning years, developers contribute across projects over time, and tools and practices are adopted in staggered fashion---precisely the structure that DiD, panel fixed effects, and synthetic control designs exploit.
Second, SE has the data infrastructure---version control systems, public repositories, platform event streams---that researchers in economics, psychology, or epidemiology often cannot access.
Finally, SE ecosystems generate rich \emph{quasi-experimental variation}: Platform policies (GitHub Actions rollout, npm security mandates, language deprecation announcements), organizational mandates, and ecosystem-wide shocks create exogenous variation that resembles natural experiments.
The field is rich in settings where design-based identification is feasible; it has simply not been exploited at scale.

The causal inference framework also creates a productive bridge between qualitative and quantitative SE research.
Every step of a causal inquiry---framing the question, proposing a DAG, choosing covariates, interpreting results---requires \emph{domain knowledge} that cannot be derived from data alone (Section~\ref{sec:stance-prior-research}).
Interviews, case studies, practitioner postmortems, and gray literature become the theory-generating engine that proposes causal structures; rigorous quantitative designs then test them; and the findings refine theory for the next iteration.
Far from marginalizing qualitative research, the toolkit makes its role in theory-building explicit and essential.

\subsection{Limitations and Open Frontiers}

Raising empirical standards cannot prevent the superficial or ritualistic adoption of methods---applying DiD without understanding parallel trends, or drawing a DAG without genuine engagement with the causal structure.
This concern is not hypothetical: The SE community has experienced analogous problems with the adoption of qualitative methods, where the label is adopted but the disciplined reasoning behind the method is not~\cite{DBLP:journals/tosem/RobillardAEGLNNSSS24}.
The antidote is peer review informed by the empirical standard (Appendix~\ref{sec:appendix-standard}) and a culture of substantive engagement with identification assumptions---evaluating whether a study's parallel trends argument is plausible, not merely whether the words ``parallel trends'' appear.

Several SE-specific challenges deserve particular attention as open research frontiers.
First, SE treatments are often \emph{compound}: ``Adopting AI coding tools'' bundles the tools themselves, the workflow reorganization, and the team characteristics that drive adoption (Section~\ref{sec:study-setting}).
Developing decomposition strategies for disentangling bundled interventions in SE settings is an important methodological direction.
Second, interconnected software ecosystems create \emph{SUTVA violations}: Shared libraries propagate effects across projects, developer mobility carries practices between teams, and upstream changes spill over into downstream dependencies.
Adapting causal inference methods under network interference~\cite{savje2021average, forastiere2021identification} to SE's sociotechnical dependency structures is an open frontier where SE-specific methodological contributions could flow back to the broader causal inference community.
Third, our worked example reports a single ATT, but the effect of AI coding tools is almost certainly \emph{heterogeneous}---varying across programming languages, team sizes, project maturity, and adoption cohorts.
Machine-learning estimators for conditional average treatment effects---causal forests~\cite{wager2018estimation} and double/debiased machine learning~\cite{chernozhukov2018double}---offer tools for characterizing this heterogeneity, but they inherit the identifying assumptions of simpler estimators and are therefore most useful once identification has been established.
Systematically mapping effect heterogeneity in SE settings---where treatment effects may interact with ecosystem position, developer experience, or project governance structures---is a natural next step for studies that have cleared the identification bar.
Finally, even as the community improves identification, corresponding investment in measurement validity---construct definition, proxy validation, and triangulation across multiple outcome measures~\cite{lawlor2016triangulation}---is essential.
A perfectly identified design with a bad outcome proxy still produces misleading conclusions.
Identification and measurement are orthogonal challenges (Section~\ref{sec:primer-stance}); this tutorial addresses the former, but the latter demands equal attention.

\section{Conclusion}

No single observational study resolves a causal question definitively.
The field builds general causal knowledge through \emph{triangulation}: accumulating credible local estimates across complementary designs, data sources, and populations~\cite{shadish2002experimental}---just as economics built its evidence base on minimum wage effects and returns to education through dozens of independent, credible studies that converged~\cite{angrist2010credibility}.
SE's diversity of ecosystems---different languages, platforms, team structures, development cultures---makes it well suited to this strategy: The same causal question can be studied in many independent settings, each with its own identification strategy, and convergence across them is far more persuasive than any single study claiming a universal effect.

The causal inference toolkit's greatest contribution to SE is not a specific method but a way of thinking: making identifying assumptions \emph{transparent and open to scrutiny} so that the strength of evidence is clearly communicated and the field can iteratively develop stronger designs.
We envision this adoption bringing productive advancements on a broad range of important intervention-outcome questions---from the effect of AI coding tools to the impact of development practices, governance structures, and ecosystem policies---moving the field from the blanket disclaimer that ``correlation does not imply causation'' toward the constructive discipline of asking: \emph{Under what assumptions does this evidence support a causal interpretation, and how can the next study relax those assumptions?}
